\centerline{\bf \Huge Теория Чисел. Продолжение}

\begin{flushright}
        \scriptsize{Единственный счастливый предмет, который вам понадобится, уже находится внутри — УВЕРЕННОСТЬ\\ Мидорима Синтаро}
\end{flushright}

\problem{\textbf{1(1)}} Найдите два таких числа, чтобы при умножении одного из них на 12, а другого на 16 получились равные произведения.

\problem{\textbf{2(1)}} Найдите наименьшее натуральное $N$
такое, что $N$ делится (без остатка) на 12, $N$ + 2 — на 14, $N$ + 4 — на 16, а $N$ + 6 — на 18.

\problem{\textbf{3(2)}} B ряд лежат 1000 конфет. Сначала Вася съел
девятую конфету слева, после чего съедал каждую седьмую конфету, двигаясь вправо. После
этого Петя съел седьмую слева из оставшихся конфет, а затем съедал каждую девятую из них,
также двигаясь вправо. Сколько конфет после этого осталось?

\problem{\textbf{4(2)}} У Незнайки есть пять карточек с цифрами: 1 , 2 ,
3 , 4 и 5 . Помогите ему составить из этих карточек два числа — трёхзначное и двузначное —
так, чтобы первое число делилось на второе.

\problem{\textbf{5(2)}}  Сумма трёх различных наименьших делителей
некоторого числа $A$ равна 8. На сколько нулей может оканчиваться число $A$?

\problem{\textbf{6(2)}} Найдите $\tau(661500)$ и $\sigma(6561). $

\problem{\textbf{7(2)}} Использовав каждую из цифр от 0 до 9 ровно по
разу, запишите 5 ненулевых чисел так, чтобы каждое делилось на предыдущее.

\problem{\textbf{8(2)}} Вася набирает электронное письмо другу, нажимая на клавиши
пальцами одной руки в следующем порядке: большим, указательным, средним, безымянным,
мизинцем, безымянным, средним, указательным, большим, указательным, далее по циклу. Каким пальцем был набран 2012-й символ?

\problem{\textbf{9(2)}} Найдите наибольшее трёхзначное число, которое
кратно сумме своих цифр и в котором первая цифра совпадает с третьей, но не совпадает со
второй.
